\section{Development History}
\label{sec:development-history}

This appendix records development notes that were posted over time.

\subsection{Release Notes}

\subsubsection{2026-9 News --- Version 0.543 (Upcoming)}

Version \code{0.543} improves repeated headless rendering and continues work on
the fused GPU reference-orbit implementation.

\begin{itemize}
\item \textbf{Persistent CLI server.} \code{FractalSharkCli} can now run as a
  server and accept repeated render requests from additional invocations of the
  same executable. The server retains initialized renderer state and processes
  requests in FIFO order. Local IPC uses named pipes on Windows and Unix-domain
  sockets on Linux.

\item \textbf{Asynchronous PNG output and timing.} Server renders begin PNG
  encoding in the background so encoding can overlap the next frame. Client
  output reports per-frame request time, while orderly server shutdown drains
  all outstanding encoders. A PowerShell example renders four deep-zoom scenes
  and reports total end-to-end time.

\item \textbf{GPU reference-orbit work.} The fused Ref2/v2 implementation
  receives further NTT, carry-propagation, memory-layout, and test-harness
  refinements.

\item \textbf{Build targets and platform support.} Release builds explicitly
  target compute capabilities 7.5, 8.0, 8.6, 8.9, 9.0, 10.0, and 12.0. Local
  builds default to compute capability 12.0 unless overridden. Windows and
  Linux remain fully supported at feature parity.
\end{itemize}

\subsubsection{2026-9-12 News --- Version 0.542}

\begin{itemize}
\item \textbf{CLI high-precision initialization.} The command-line renderer now
  initializes MPIR's default precision before parsing and converting locations,
  matching graphical-application initialization. This fixes flat direct-coordinate
  renders and precision failures at valid deep-zoom locations.
\end{itemize}

\subsubsection{2026-9-11 News --- Version 0.541}

\begin{itemize}
\item \textbf{CUDA architecture compatibility.} CUDA code generation now
  includes a compute-capability 7.5 target, restoring the fallback path used by
  RTX 3xxx cards and other supported GPUs without a matching native image.
\end{itemize}

\subsubsection{2026-9-7 News --- Version 0.54}

Version \code{0.54} replaces the original multi-kernel GPU high-precision
arithmetic path with the fused Ref2/v2 reference-orbit implementation.

\begin{itemize}
\item \textbf{Fused GPU reference orbit.} High-precision recurrence, NTT
  multiplication, signed carry propagation, normalization, and output run in a
  persistent cooperative CUDA kernel. The superseded standalone GPU add and
  multiply implementation is removed.

\item \textbf{GPU Newton--Raphson.} Feature-finder refinement uses the same
  fused high-precision machinery, including derivative accumulation and
  checkpoint validation against the CPU reference implementation.

\item \textbf{Extreme built-in views.} Views~\#33 and~\#34 exercise the new
  implementation at extreme depth. View~\#33's coordinates are corrected, and
  View~\#34 documents a feature at approximately $10^{650452}$ magnification.

\item \textbf{Validation and tooling.} GPU arithmetic tests, profiling tools,
  checksum diagnostics, and the technical notes are expanded around Ref2/v2.
\end{itemize}

\subsubsection{2026-7-4 News --- Version 0.532}

\begin{itemize}
\item \textbf{Rendering pipeline.} The render-thread pool and frame-publication
  path are refined, and parallel PNG saving is improved so output encoding can
  proceed without unnecessarily blocking subsequent rendering.

\item \textbf{Numeric performance.} HDR scaling operations receive a bitwise
  implementation that substantially reduces per-pixel overhead.

\item \textbf{Interaction fixes.} Feature autozoom is simplified around the
  stateless zoom-to-feature workflow. Palette rotation, context rendering,
  checkpoint metadata, OpenGL presentation, and startup behavior receive
  correctness fixes.

\item \textbf{Documentation.} The technical notes are reorganized into their
  current numbered sections and appendices.
\end{itemize}

\subsubsection{2026-6-27 News --- Version 0.531}

\begin{itemize}
\item \textbf{Shared graphical behavior.} Portable command handling, file
  operations, help content, saved locations, and menu-state logic are shared by
  the Win32 and Linux front ends. The Linux application is substantially
  simplified around the same contracts and reaches feature parity with Win32.

\item \textbf{Correctness fixes.} This release fixes a GPU normalization
  regression introduced between versions~0.51 and~0.52, a heap allocator false
  overrun report, repaint-state handling, and several GUI and numeric edge cases.

\item \textbf{Release packaging.} Continuous-integration artifacts and release
  packages are expanded, and project naming is made consistent across the two
  platforms.
\end{itemize}

\subsubsection{2026-6 News --- Version 0.53}

Version \code{0.53} introduced native Linux support.

\begin{itemize}
\item \textbf{Native Linux build.} The CMake + Clang build now covers the core
  numeric library, CUDA rendering, GPU high-precision arithmetic, tests, the
  command-line renderer, and the native graphical application. Debug and
  Release configurations are built in parallel with the Visual~Studio
  solution, and Linux release binaries are targeted at Ubuntu.

\item \textbf{Native Linux GUI.} \code{FractalSharkGuiLinux} now uses Xlib and
  GLX for its window and OpenGL presentation, a native hierarchical X11 popup
  menu, and Dear~ImGui overlays for help and application dialogs. Navigation,
  drag and wheel zoom, automatic zoom, feature finding, fullscreen transitions,
  clipboard output, and the principal save/load workflows are connected to the
  shared renderer and command catalog. The Linux and Win32 applications are
  fully supported and maintained at feature parity. Menu enablement, check
  marks, and radio selections are evaluated by the same portable implementation
  on both platforms.

\item \textbf{Portable platform services.} Operating-system dependencies are
  consolidated in \code{FractalSharkPlatform}, including allocation, files,
  process and memory queries, threading, input, diagnostics, wait cursors,
  embedded resources, and native OpenGL contexts. Shared library interfaces use
  portable window and geometry types instead of exposing Win32 types.

\item \textbf{Linux validation and packaging.} Linux CI builds Debug and
  Release, runs the portable CPU unit tests, smoke-tests the headless renderer,
  checks the release executables for unresolved dependencies, and packages
  \code{FractalSharkCli} with \code{FractalSharkGuiLinux}. CUDA-specific tests
  are compiled but remain a hardware validation step because hosted runners do
  not provide a physical GPU.
\end{itemize}

\subsubsection{2026-5-3 News --- Version 0.52}

Version \code{0.52} focuses on GPU-accelerated Newton--Raphson and begins the
cross-platform port to Linux.

\begin{itemize}
\item \textbf{GPU Newton--Raphson.}  A full GPU implementation of Newton's method
  for the feature finder is now integrated end-to-end: reference-orbit
  computation via NTT, first- and second-derivative accumulation, and
  convergence testing all run on the GPU.  An occupancy calculation bug is
  fixed, and a CPU-side ``reference reference orbit'' validates correctness
  against MPIR ground truth.

\item \textbf{Newton--Raphson checkpointing.}  The NR inner loop now supports
  save/resume with full-precision derivative state.  CPU and GPU checkpoint
  formats are compatible, so an interrupted computation can be resumed on
  either backend.

\item \textbf{Linux port (in progress).} A CMake + Clang build system is
  introduced alongside the existing Visual~Studio solution. At the time of
  this release, a GTK-based GUI was proposed and the port was incomplete.
  Version~0.53 superseded that proposal with the native Xlib/GLX and
  Dear~ImGui application described above. CPU-only unit tests already passed on
  Linux at this stage.

\item \textbf{Miscellaneous.}  Extracted \code{FeatureFinderOrchestrator} from
  the main class; improved the heap allocator panic path; parallelized the
  CI build; fixed Imagina file loading and RDP rendering regressions.
\end{itemize}

\subsubsection{2026-2-14 News --- Version 0.51}

Version \code{0.51} introduces the \textbf{Feature Finder} (\cref{sec:feature-finder}), a major new tool
for exploring the Mandelbrot set, along with several usability improvements.

\begin{itemize}
\item \textbf{Feature Finder.}  Move mouse near a feature, select feature
  finder, and \FractalShark{} will locate the nearest interesting feature---a
  mini-Mandelbrot or satellite bulb---using Newton's method (\cref{subsec:ff-newton}).  A line is drawn
  from the cursor to the detected feature, making its location visible.
  The finder works at any zoom depth.  Currently, the user manually chooses between
  direct iteration, perturbation theory, and linear approximation depending on
  the current zoom depth.  In the future, this choice will be automatic.

\item \textbf{Scan mode.}  Instead of a single click, scan mode surveys a grid
  of points across the screen and marks up to 144 discovered features with small
  circles.  This provides a quick overview of what is nearby before deciding
  where to zoom next.

\item \textbf{Auto-zoom to feature.}  Once a feature has been found the user can direct
  the program to zoom into it automatically.  The iteration count is ramped up
  smoothly during the zoom so the image stays sharp as the view approaches the target.
  This feature remains hacky and experimental but if it works it's fun to
  play with.

\item \textbf{Mouse-wheel zoom.}  Scroll the mouse wheel forward to zoom in
  toward the cursor, or backward to zoom out.  The code is rewritten.

\item \textbf{Rendering improvements.}  Parts of the OpenGL display layer are
  rewritten for better reliability and to support the new feature-finder
  overlays (lines and circles drawn on top of the fractal).

\item \textbf{Under the hood.}  Extended the derivative math in the
  linear-approximation pipeline so the feature finder works at extreme depths.
  Improved precision tracking during zoom.  Added a \code{.clang-format}
  configuration for consistent code formatting.
\end{itemize}

\subsubsection{2026-1-18 News --- Version 0.501}

\begin{itemize}
\item \textbf{Basic Wine compatibility.}  Tested without a GPU and working under
Wine 9.0 on Linux.   There are still some quirks, but basic rendering works in
CPU-only mode. Whether \FractalShark{} can work with Wine + CUDA remains unknown as
an appropriate test platform is not currently available.
\item \textbf{Proper fallback when no GPU is present.}  The program now detects lack of a CUDA
device and falls back to CPU-only rendering, with a warning message.
\item \textbf{Fixed some windowing bugs.}
\item \textbf{Improved these notes.}
\item \textbf{Linked at FractalForums.}  See \url{https://fractalforums.org/index.php?topic=5564.0}
\end{itemize}

\subsubsection{2025-12-29 News --- Version 0.5}
\begin{itemize}
  \item \textbf{GPU reference-orbit prototype.} Version~0.5 introduced the v1
  CUDA reference-orbit implementation, which used separate high-precision
  arithmetic phases.  That implementation has since been removed and
  superseded by the fused Ref2/v2 implementation documented in
  \cref{subsec:ref-orbit-gpu}.

  \item \textbf{Project plumbing.} Hooks up GitHub Actions to produce “official”
  builds, plus some refactoring; notes a startup windowing quirk.
\end{itemize}

\subsubsection{2025-07-28 News --- Version 0.46}

Version \code{0.46} was published as a pre-release, focusing on refactoring,
reference compression experimentation, save-file compatibility, and test infrastructure.

\begin{itemize}

  \item Continued internal \textbf{refactoring} of core rendering and control logic.
  
  \item Experimental support for \textbf{Imagina-style “max” reference compression},
        as an alternative to the existing “simple” scheme.
  
  \item Added preliminary support for \textbf{Imagina-compatible save files}.

  \item Enhanced internal \textbf{test infrastructure} with broader coverage and tooling.
    
  \item Work in progress toward additional infrastructure improvements.

\end{itemize}

This release was published on GitHub on July 28, 2025 (UTC) and is tagged as a
pre-release build.  


\subsubsection{2024-04-20 --- Version 0.45}

\begin{itemize}
  
  \item Adds a first cut at reference compression for an intermediate “perturbed
  perturbation” reference-orbit strategy.

  \item Includes three internal variants (v1: uncompressed/3 threads; v2:
  uncompressed/4 threads; v3: compressed/4 threads).

  \item v3 becomes the default at very deep “auto” perturbation settings;
  guidance given for when to stick to the older MT2 + periodicity workflow.

\end{itemize}


\subsubsection{2024-03-31 News --- Version 0.44}

Version \code{0.44} represents a substantial internal refactor emphasizing memory management, allocator control, and reference-orbit performance.

\begin{itemize}

  \item Removed the Boost dependency and introduced a custom
  \textbf{HighPrecision} wrapper.

  \item Improved reference-orbit performance, especially for low-precision,
  high-period locations.

  \item Improved performance of the \textbf{perturbed perturbation} algorithm.

  \item Introduced a custom file-backed \textbf{bump allocator} (\cref{sec:disk_backed_growable_vectors}) with stable
  interior pointers.

  \item Extended \code{GrowableVector} to support allocator-style usage
  without increasing committed memory.
  
  \item Updated reference orbit and linear approximation save-file formats.

  \item Added an \textbf{automatic perturbation mode} switching between single-threaded and multithreaded execution.

  \item Instrumented for memory-leak detection and fixed all known leaks.

  \item Fixed a long-standing correctness bug in BLAv1 dating back to version
  0.21.

  \item \textbf{Clarified ``perturbed perturbation'' design.} The implementation
  is framed as storing an intermediate-resolution reference orbit (e.g.\
  $\sim$500 bits), then using low-precision perturbation off that intermediate
  orbit to generate subsequent nearby reference orbits more efficiently.

\end{itemize}


\subsubsection{2024-03-03 News --- Version 0.43}
\begin{itemize}
  
  \item Significant performance wins for lower-precision, high-period
  reference-orbit computation.
  
  \item Adds a custom per-thread bump allocator for MPIR high-precision numbers.
  
  \item Rebalances work across the existing 3-thread multithreaded
  reference-orbit approach.

  \item \textbf{Allocator usage guidance.} The bump allocator is primarily
  beneficial for low-precision, high-period reference orbits; at very high
  precision the arithmetic dominates and allocation overhead is negligible. A
  suggested heuristic is to use the fixed-block bump allocator below roughly
  1000 digits and fall back to the global allocator above that.

  \item \textbf{Empirical sizing recommendation.} Tested successfully on a $\sim
  1\mathrm{e}{6000}$ location with about 4\,MB of bump-allocator space;
  recommendation is to use the smallest practical buffer (cache behavior matters).

  \item \textbf{Implementation refinement note.} In the current multithreaded
  orbit implementation, most allocations may occur on the main thread;
  thread-local allocation may be unnecessary. A few small allocator issues were
  identified for cleanup in the following release.

\end{itemize}


\subsubsection{2024-02-24 News --- Version 0.42}

Version \code{0.42} introduced major improvements to linear approximation performance, benchmarking, and configurability.

\begin{itemize}
  \item Added \textbf{multithreaded linear approximation table generation}.
  \item Significantly improved performance for high-period locations, especially with reference compression enabled.
  \item Added fine-grained benchmarking and performance breakdowns.
  \item Enabled regeneration of linear approximation tables independently of per-pixel data.
  \item Introduced adjustable linear approximation presets trading memory, accuracy, and performance.
\end{itemize}

This release corrected a prior precision regression and restored high-accuracy defaults.


\subsubsection{2024-01-15 News --- Version 0.41}

\begin{itemize}
  \item Fixes a minor bug in reference compression / GPU.
  \item Memory-use overhaul for storing LA tables and reference orbit to minimize RAM needed for hard views (e.g., View \#27).
  \item Default mode now uses temporary files + memory-mapped IO; avoids needing 2$\times$ memory during resize/copy and reduces committed-memory requirements with no measurable overhead per the release note.
\end{itemize}


\subsubsection{2024-01-01 News --- Version 0.4}

\begin{itemize}
  \item First working implementation of \textbf{runtime} reference compression
  (decompress-on-demand during rendering).

  \item Motivated as reducing end-to-end RAM requirements substantially on very
  high-period locations; not enabled by default for “Auto” rendering and
  described as work-in-progress.

  \item \textbf{Motivation: memory pressure.} Reference compression is
  explicitly framed as a response to high RAM requirements at difficult views
  (notably after enabling 64-bit iteration counts).
\end{itemize}


\subsubsection{2023-12 News --- Version 0.32}

Version \code{0.32} focused on usability, correctness, and incremental architectural cleanup.

\begin{itemize}
  \item Added \textbf{progressive rendering}, allowing partial image updates
  during long GPU renders.
  \item Added a \textbf{2x32 non-HDR linear approximation} path for numerically
  difficult shallow zooms.
  \item Fixed a subtle but severe correctness bug in the 2x32 implementation
  present in version 0.31.
  \item Added a basic regression test suite.
  \item Refactored code and restored buildability across inactive code paths.
\end{itemize}


\subsubsection{2023-12-09 News --- Version 0.31}

Version \code{0.31} expanded the set of available rendering algorithms and
substantially improved default performance at shallow and mid-depth zooms.

\begin{itemize}
  \item Added \textbf{1x32 linear approximation / perturbation} for
  high-performance shallow zooms.
  
  \item Added \textbf{1x64 linear approximation / perturbation}, primarily for
  correctness testing.
  
  \item Implemented \textbf{automatic render-algorithm selection} based on zoom
  depth.
  
  \item Eliminated unnecessary reference-orbit copying, significantly improving
  frame-to-frame performance.
  
  \item Fixed several bugs, including 2x32 correctness issues and unaligned
  HDRx64 memory access.

  \item \textbf{Auto-selection policy detail.} Auto mode is described as staging
  from direct 32-bit rendering at very shallow zooms, to 1x32 perturb-only, to
  1x32 LA, and finally 1x32 float-exp LA at deeper zooms, with the goal of
  improving default performance below roughly $10^{38}$.

  \item \textbf{Reference-orbit reuse optimization.} Removing unnecessary
  reference-orbit copying is called out as a major win for deep-zoom
  frame-to-frame responsiveness when reusing a reference orbit.

\end{itemize}


\subsubsection{2023-11-26 News --- Version 0.3}

\begin{itemize}
  
  \item Adds HDRx2x32 linear approximation: faster than LA/HDRx64 on consumer
  cards while retaining nearly native-64-bit precision.
  
  \item Notes that a “Debug View 20” anomaly is resolved under HDRx2x32,
  supporting the hypothesis that HDRx32 precision limits caused the issue.
  
  \item Misc performance improvements and bug fixes elsewhere.

  \item \textbf{Debug View 20 consistency improvement.} The ``Debug View 20''
  anomaly is reported as resolved under HDRx2x32, supporting the conclusion that
  the earlier HDRx32 behavior was a precision-limit issue.

  \item \textbf{HDRx2x32 precision characterization.} Described as ``Linear
  Approximation + Float EXP'' using a pair of 32-bit floats plus an exponent,
  with an estimated combined precision of roughly 46 bits.

\end{itemize}


\subsubsection{2023-11-12 News --- Version 0.24}

\begin{itemize}
\item \textbf{UI/UX improvements.} Adds a hotkey help view and general UI
cleanup.
\item \textbf{Random palette generation.} Adds a random palette generator; notes
the palette system remains intentionally constrained and may be expanded.
\item \textbf{Rendering diagnostics.} Adds a ``Show Rendering Details'' view for
inspecting current algorithm/settings.
\item \textbf{Early HDRx2x32 work.} Begins a GPU HDR 2x32 float implementation
(explicitly noted as not working yet).
\item \textbf{Crash diagnostics.} Automatically generates a crash dump on
unhandled exceptions.
\end{itemize}


\subsubsection{2023-11-05 News --- Version 0.23}
\begin{itemize}
  \item Adds switchable 64-bit iteration count support (enabling tens of
  billions of iterations per pixel, leveraging LA).
  \item Moves antialiasing/coloring into its own CUDA kernel for better
  performance.
  \item Multiple memory-reduction optimizations (systems-level) and related perf
  work.
  \item Adds reference-orbit save/load (no compression yet), useful for
  debugging deep spots.
  \item Notes a known bug manifesting in built-in View \#20 under a particular
  configuration (64-bit iterations + default HDRx32 LA~v2), with long render
  times due to slow reference-orbit calculation.
\end{itemize}


\subsubsection{2023-10-14 News --- Version 0.22}

\begin{itemize}

  \item Significant LA~v2 performance improvements (bug fixes); notes BLAv1 still
  winning in some cases.

  \item Fixes an LA~v2 correctness bug; release note claims no known correctness
  bugs remaining.

  \item \textbf{Performance data point.} Reports a minibrot example improving
  from roughly 470\,ms to 370\,ms (render-only; reference orbit excluded) on the
  author’s machine.

\end{itemize}


\subsubsection{2023-10-08 News --- Version 0.21}
\begin{itemize}
  \item Performance improvements to BLA and LA~v2.
  \item Memory-usage improvements (commit size / virtual address space issues
  called out).
  \item Generates more CUDA kernels for older GPUs (system requirements
  unchanged).
  \item Notes LA~v2 is strong for many deep-zoom cases but still behind older BLA
  in some shallower cases.
\end{itemize}


\subsubsection{2023-09-16 News --- Version 0.2}

Version \code{0.2} introduced a major algorithmic milestone: the integration
of Imagina's linear approximation implementation as a CUDA kernel.

\begin{itemize}
  \item Added Imagina-based linear approximation (\code{LAv2}), now the
  default rendering algorithm.
  \item Resolved known correctness bugs in the new approximation path.
  \item Fixed a major performance regression caused by an earlier implementation oversight.
  \item Achieved large performance gains at deeper zoom levels compared to the
  existing bilinear approximation.
\end{itemize}

While the legacy bilinear approximation remains faster at very shallow zooms,
the new approach dominates once meaningful magnification is reached.


\subsubsection{2023-07-22 News --- Version 0.11}

Version \code{0.11} was released with a focus on performance improvements and
improved diagnostics.

\begin{itemize}
  \item Approximately \textbf{+10\% performance improvement} in the HDRx32 CUDA
  rendering path.
  \item Fixed several minor bugs affecting correctness and stability.
  \item Improved CUDA error reporting by displaying descriptive error strings
  instead of numeric codes.
\end{itemize}

A long-duration technology demonstration video accompanied this release,
showcasing a zoom to approximately $10^{4000}$ using the default HDRx32/BLA CUDA
kernel.  Rendering and post-processing each required roughly one and a half
days.
